\chapter{Considerações Finais}\label{cap:Considerações}

%\section{Conclusões}

Embora este estudo tenha avançado na compreensão das interações entre as ondas atmosféricas e a contribuição estratosférica para eventos de seca, existem vários aspectos que merecem investigação adicional.

Primeiramente, o acoplamento troposfera-estratosfera, especialmente sob cenários de mudança climática, ainda não foi completamente explorado. Estudos futuros devem focar nesse acoplamento, analisando como os processos estratosféricos influenciam a dinâmica atmosférica na troposfera.

Além disso, a dinâmica não linear entre os modos normais e a interação entre ondas de diferentes escalas, como as de pequena e média escala, também deve ser investigada. A análise não linear pode fornecer novas explicações para os mecanismos que provocaram a estacionariedade das ondas em grande escala. As mudanças da velocidade de propagação das ondas certamente são influenciadas pela presença de forçantes (por ex., a liberação de calor latente). Entretanto, a contribuição dos termos não lineares também pode ter um papel fundamental, incluindo o processo de quebra de ondas, interação não linear entre ondas \citep{Enver2017}, promovendo mudanças abruptas no clima ou estacionaridade.

Outro aspecto a ser explorado é a baixa energia observada em alguns modos verticais elevados, como no caso do modo \(m = 33\), cuja estrutura vertical concentra a amplitude no topo (\(\sigma = 0\)). Esse comportamento sugere que se trate de um modo computacional gerado pela discretização da equação da estrutura vertical.

O fato de aparecerem mudanças muito significativas no comportamento de alguns modos normais nas diferentes fases do episódio da seca de 2013/14 precisa ser avaliado em outros episódios que levaram a períodos extremos na precipitação no SE do Brasil e a possível conexão global.

Este trabalho concentrou-se em uma análise diagnóstica dos processos que culminaram no fenômeno extremo. Os resultados obtidos indicam um potencial aspecto preditivo para a ocorrência de transições no regime climático. Em outras palavras, o monitoramento da evolução dos padrões atmosféricos no espaço de fase (ou seja, no espaço formado pelos coeficientes de expansão dos modos normais) pode, por meio de modelos simplificados ou do uso de Inteligência Artificial, fornecer indicativos das transições com certa antecedência. Nesse contexto, também é fundamental investigar o comportamento das transições do estado atmosférico sob a perspectiva dos modos normais, especialmente em modelos globais de previsão. Surge, então, a questão: será que os modelos de previsão climática conseguem reproduzir as trocas de energia entre ondas de Rossby, ondas de gravidade inerciais, ondas Kelvin e ondas mistas Rossby-Gravidade observadas?
